\section{Sampled source factors and participant-pair labels}
\label{sec:peps_sampled_source_network}

This section records the sampled source factorization and the finite labels
used in the network construction of
\cite[proof of Theorem~5.2]{OpenAI2026PolynomialPEPS},
\texttt{04-compression.tex}, lines 233--251, 279--309 and 565--588.
Let $\mathcal G$ be the nonprivate gate occurrences of a source-only circuit.
For $g\in\mathcal G$, write $P_g$ for its participating set, $I_g$ for its
branch set and $n_g$ for its number of common source slots. The source
occurrences form $\mathcal E=\bigsqcup_{g\in\mathcal G}\{1,\ldots,n_g\}$.
For $e=(g,i)$, denote its ordered endpoint parties by $(p_e,q_e)$.

\begin{definition}[Two sampled source factors]
    \label{def:peps_sampled_source_factors}
    \lean{TNLean.PEPS.PairEffect.SourceCircuit.sampledEndpointLeft}
    \lean{TNLean.PEPS.PairEffect.SourceCircuit.sampledEndpointRight}
    \leanok
    Fix $k\ge0$, one Gaussian array $g^{e,\xi,\zeta}_{j,a,c}$ at each source
    occurrence and ordered pair $(\xi,\zeta)\in I_g\times I_g$, and endpoint
    coordinate matrices $E_{e,\xi},F_{e,\xi}$ with real nonnegative weights
    $\lambda_{e,\xi,a}$. For $j\in\{1,\ldots,k\}$ define the rectangular
    matrices $A^{e,\xi,\zeta}_j$ and $B^{e,\xi,\zeta}_j$ by
    \begin{align}
        A^{e,\xi,\zeta}_j(a,c)
        &= (\lambda_{e,\xi,a}\lambda_{e,\zeta,c})^{1/4}
           g^{e,\xi,\zeta}_{j,a,c}, \\
        B^{e,\xi,\zeta}_j(a,c)
        &= (\lambda_{e,\xi,a}\lambda_{e,\zeta,c})^{1/4}
           \overline{g^{e,\xi,\zeta}_{j,a,c}}.
    \end{align}
    Their actual endpoint factors are
    \begin{align}
        U^{e,\xi,\zeta}_j
        &=E_{e,\xi}A^{e,\xi,\zeta}_jE_{e,\zeta}^{*}, \\
        V^{e,\xi,\zeta}_j
        &=F_{e,\xi}B^{e,\xi,\zeta}_jF_{e,\zeta}^{*}.
    \end{align}
    Both factors use the same array, the same ordered branch pair and the
    same sample index. Their definitions make sense for each fixed array;
    no probabilistic assertion is needed here.
\end{definition}

\begin{theorem}[Actual sampled source factorization]
    \label{thm:peps_sampled_source_factors}
    \lean{TNLean.PEPS.PairEffect.SourceCircuit.sampledSourceMatrix_eq_endpoint_average}
    \leanok
    \uses{def:peps_sampled_source_factors}
    The sampled source matrix at $e$ and $(\xi,\zeta)$ is
    \begin{align}
        \widehat X^{e,\xi,\zeta}
        =k^{-1}\sum_{j=1}^{k}
          U^{e,\xi,\zeta}_j\otimes V^{e,\xi,\zeta}_j.
        \label{eq:peps_sampled_source_factors}
    \end{align}
    At $k=0$, both the empty sum and the inverse of zero are taken to be
    zero, so the identity also covers that case.
\end{theorem}

\begin{proof}\leanok
    \uses{def:peps_sampled_source_factors}
    Substitute the two endpoint factors into the definition of the sampled
    source matrix.
\end{proof}

\begin{theorem}[Actual source slots and pair-sample edges]
    \label{thm:peps_sampled_source_pairs}
    \lean{TNLean.PEPS.PairEffect.SourceCircuit.sourceSlotPair_bijective}
    \lean{TNLean.PEPS.PairEffect.SourceCircuit.sourcePairAt_injective}
    \lean{TNLean.PEPS.PairEffect.SourceCircuit.sourceSlotPairEquiv}
    \lean{TNLean.PEPS.PairEffect.SourceCircuit.sourceSlotPairEquiv_endpoints}
    \lean{TNLean.PEPS.PairEffect.SourceCircuit.sourceSampleEquiv}
    \lean{TNLean.PEPS.PairEffect.SourceCircuit.sourceSampleEquiv_gate}
    \lean{TNLean.PEPS.PairEffect.SourceCircuit.sourceSampleEquiv_parties}
    \leanok
    \uses{def:peps_source_slots,def:peps_source_circuit_positions}
    Let $\binom{P_g}{2}$ denote the unordered pairs of distinct members of
    $P_g$. The map $i\mapsto\{p_{(g,i)},q_{(g,i)}\}$ is a bijection from
    the common source slots of $g$ to $\binom{P_g}{2}$. It gives the
    bijection
    \begin{align}
        \Phi:\mathcal E&\longrightarrow
           \bigsqcup_{g\in\mathcal G}\binom{P_g}{2}, &
        (g,i)&\longmapsto(g,\{p_{(g,i)},q_{(g,i)}\}).
        \label{eq:peps_sampled_source_pairs}
    \end{align}
    Thus every pair-sample edge retains the original gate occurrence and
    exactly the two actual source endpoint parties. Empty participating
    sets have no common source slots and no pair-sample edges.
\end{theorem}

\begin{proof}\leanok
    \uses{thm:peps_source_only_resources}
    Induct on the circuit composition. At a prepared gate, the common
    slot pairs have no repetitions. The injective inclusion of its local
    parties into the original party set preserves this property. There
    are $\binom{|P_g|}{2}$ slots and the same number of unordered distinct
    participant pairs, so the injective map is bijective. Composition
    takes disjoint unions of these bijections, and spectator registers
    leave them unchanged.
\end{proof}

\begin{theorem}[Preservation of the shared sample index]
    \label{thm:peps_sampled_source_indices}
    \lean{TNLean.PEPS.PairEffect.SourceCircuit.sourceSampleIndices_apply}
    \lean{TNLean.PEPS.PairEffect.SourceCircuit.sourceSampleIndices}
    \leanok
    \uses{thm:peps_sampled_source_pairs}
    For an assignment $j:\mathcal E\to\{1,\ldots,k\}$, define the
    edge assignment by $\widetilde j=j\circ\Phi^{-1}$. Then
    $\widetilde j(\Phi(e))=j(e)$ at every actual source occurrence $e$.
\end{theorem}

\begin{proof}\leanok
    \uses{thm:peps_sampled_source_pairs}
    The inverse identity $\Phi^{-1}(\Phi(e))=e$ gives the result.
\end{proof}

\begin{theorem}[Dimensions of the constructed virtual labels]
    \label{thm:peps_sampled_source_dimensions}
    \lean{TNLean.PEPS.PairEffect.SourceCircuit.card_densityBranchLabels}
    \lean{TNLean.PEPS.PairEffect.SourceCircuit.card_densityBranchLabels_le}
    \lean{TNLean.PEPS.PairEffect.SourceCircuit.card_sourceNetworkAlphabet_le}
    \leanok
    At a gate $g$, the common density label $(\xi,\zeta)$ has
    $|I_g|^2$ values. Share this label along its star links and share an
    index in $\{1,\ldots,k\}$ along each pair-sample edge. If $|I_g|\le M$
    on the selected gates, every resulting virtual alphabet has
    cardinality at most $\max\{M^2,k\}$. This bound is independent of
    private-memory dimensions and source Schmidt ranks.
\end{theorem}

\begin{proof}\leanok
    The ordered branch-pair set is $I_g\times I_g$, of cardinality
    $|I_g|^2\le M^2$. A pair-sample alphabet has exactly $k$ values.
    Taking the maximum gives the stated bound.
\end{proof}

\begin{theorem}[Scalar contribution of a gate with no participating parties]
    \label{thm:peps_empty_party_gate_scalar}
    \lean{TNLean.PEPS.PairEffect.SourceCircuit.eval_gate_heq_coeffSum_smul_id_of_participants_eq_empty}
    \lean{TNLean.PEPS.PairEffect.PreparedSourceGate.eval_heq_coeffSum_smul_id_of_isEmpty}
    \lean{TNLean.PEPS.PairEffect.PreparedSourceGate.norm_coeffSum_le_one_of_isEmpty}
    \lean{TNLean.PEPS.PairEffect.SourceCircuit.eval_gate_eq_coeffSum_smul_id_of_isEmpty}
    \leanok
    Consider a prepared gate $G=\sum_{\xi\in I}c_\xi M_\xi$ on a finite
    local party set $C$, placed in the original party set by an injective
    inclusion $\iota:C\to P$. Suppose its actual participating set
    $\iota(C)$ is empty. Put $s_G=\sum_{\xi\in I}c_\xi$. Its local input
    and output memories are the scalar space $\C$, and
    \begin{align}
        G&=s_G\operatorname{Id}_{\C}, & |s_G|&\le1.
        \label{eq:peps_empty_gate_scalar}
    \end{align}
    Beside any spectator memory $\mathcal H$, the literal placed gate is
    $s_G\operatorname{Id}_{\mathcal H}$. This includes an empty branch
    set $I$, when $s_G=0$ and both operators vanish.
    These are the zero-party gates permitted by the classification in
    \cite[proof of Theorem~5.2]{OpenAI2026PolynomialPEPS},
    \texttt{04-compression.tex}, lines 24--28 and 233--251.
\end{theorem}

\begin{proof}\leanok
    If $q\in C$, then $\iota(q)$ belongs to the participating set, a
    contradiction. Thus $C$ is empty. Every register must have an owner,
    so both local layouts are empty. A source-only word on this party
    set has no local map, source preparation, exchange or framed
    register, and hence acts as $\operatorname{Id}_{\C}$. Summing the
    branch coefficients gives $G=s_G\operatorname{Id}_{\C}$. Since
    $G$ is a contraction and $\|\operatorname{Id}_{\C}\|=1$, one has
    $|s_G|=\|G\|\le1$. Tensoring with the spectator identity and applying
    the canonical owner and tensor identifications yields
    $s_G\operatorname{Id}_{\mathcal H}$.
\end{proof}

% Remaining chronological contraction identity. No formalization tag is
% attached to the following requirement.
\paragraph{Chronological contraction.}
For the original branch coefficients $c_{g,\xi}$, retain the scalar
contribution of the gates with no participating parties:
\begin{align}
    c_{\rm empty}
    =\prod_{\substack{g\in\mathcal G\\P_g=\varnothing}}
      \left(\sum_{\xi\in I_g}c_{g,\xi}\right).
    \label{eq:peps_sampled_empty_scalar}
\end{align}
To identify the sampled physical density for the whole-party normalized
product input with the operator network, one must construct each party's
local tensor $T_p(\omega)$ by contracting its input factor, chronological
local maps, private preparations, intermediate memories and final discarded
registers. The further identity required, which is not established by the
results of this section, is
\begin{align}
    \widehat\rho(\omega)
    =|c_{\rm empty}|^2\,
      \operatorname{Contr}\bigl((T_p(\omega))_{p\in P}\bigr).
    \label{eq:peps_sampled_chronology_required}
\end{align}
Here the contraction uses the common gate labels and shared sample labels
specified above. The source averaging factors $k^{-1}$ are included once
per source in the local tensors. Gates with no participating parties
contribute through (\ref{eq:peps_sampled_empty_scalar}) without introducing
star or sample links. This requirement also includes $P=\varnothing$,
with the empty tensor contraction equal to $1$.
